Lung adenocarcinoma (LUAD) develops through a defined morphological sequence from atypical
adenomatous hyperplasia (AAH) to invasive adenocarcinoma, but the transcriptional programmes
accompanying this transition have mostly been characterised in bulk tissue or in tumour
tissue alone. We profiled paired samples spanning the full sequence by single-nucleus RNA
sequencing (413,697 nuclei, 23 patients) and spatial transcriptomics (56 Visium sections, 25
patients), asking how the transcriptional state of the tissue changes as invasion proceeds
and whether that state can be reversed \emph{in silico} by known perturbations.

The atlas resolves into 39 second-level subtypes and the tissue into seven recurring
niche-domain archetypes ($K^{*}=7$), six of which are ordinary lung structures and only one
of which is restricted to invasive disease. A confound then shaped the remaining analyses:
domain identity, marker stability and the apparent stage-dependence of the spatial programmes
all scale with sequencing depth, and depth is in these sections a proxy for tissue density
rather than a purely technical artefact ($\rho=-0.89$ between a domain's depth and the
fraction of its markers surviving depth matching). Several apparently positive results
disappeared once depth was matched. The invasive domain's ECM/interstitial identity survives
that check: within patients the fibroblast ECM programme is higher in invasive disease in 20
of 22 paired cases ($P=0.0033$ against random gene sets), and the result is unchanged when
the genes added by hand are removed and under every leave-one-out variant.

Querying a patient-paired disease axis against the LINCS compound library with three scoring
schemes returned 34 compounds in common, dominated by PI3K/mTOR inhibitors, with a
split-half reproducibility of $\rho=0.962$ against a permutation null of $\pm0.045$ and an
agreement of $\rho=0.735$ between a single-cell-derived and a spatially derived axis. A
positive control, the recovery of known regulators of an official benchmark (TBXT ranked 1st
of 8,450), establishes that the pipeline recovers true positives, so the negative results
reported here are properties of the data rather than of the method. Those negatives are
stated explicitly: single-cell copy-number inference produced no usable
malignant/non-malignant separator, the cohort-level spatial copy-number difference did not
survive depth matching, a developmental trajectory ran opposite to that of the source study,
gene-level prediction returned no signal on the pooled axis and only one perturbation
survives compartment-wise gating, and a ligand-level analysis produced no usable pair.
